% =========================================================
% Proof: Inverse Function Theorem (Monotone Version) — Global Form
% Source: volume-iii/analysis/differentiation/notes/notes-inverse-function.tex
% Mode: standard
% =========================================================

\subsection*{Inverse Function Theorem (Monotone Version) — Global Form}
\label{prf:inverse-function-theorem-global}

\begin{remark*}[Return]
\hyperref[thm:inverse-function-theorem-monotone]{$\leftarrow$ Back to Theorem (Inverse Function Theorem, Monotone Version) in Notes}
\end{remark*}

\begin{proof}
Let $I \subseteq \mathbb{R}$ be an interval, let $f : I \to \mathbb{R}$ be
strictly monotone and differentiable on $I$ with $f'(x) \neq 0$ for all
$x \in I$, set $J := f(I)$, and let $g : J \to \mathbb{R}$ be the inverse
of $f$. The claim is that $g$ is differentiable on $J$ and
$g' = 1/(f' \circ g)$.

Let $d \in J$ be arbitrary. Since $g$ is the inverse of $f$, there exists a
unique $c \in I$ with $f(c) = d$, namely $c = g(d)$. By hypothesis $f$ is
differentiable at $c$ with $f'(c) \neq 0$, so by
\ByThm{Inverse Function Theorem (Monotone Version)}, $g$ is differentiable
at $d$ and
\[
  g'(d) = \frac{1}{f'(c)} = \frac{1}{f'(g(d))} = \frac{1}{(f' \circ g)(d)}.
\]
Since $d \in J$ was arbitrary, $g$ is differentiable at every point of $J$
and the derivative function satisfies $g' = 1/(f' \circ g)$ on $J$.
\end{proof}

\begin{remark*}[Proof shape]
The global statement is obtained by instantiating the pointwise theorem at
an arbitrary $d \in J$ and substituting $c = g(d)$. No new technique is
involved — the proof is a quantifier argument over the pointwise result.

\FlashProof{1. Let $d \in J$ be arbitrary; set $c = g(d)$, the unique
  preimage.
  2. Apply the pointwise Inverse Function Theorem at $c$:
  $g'(d) = 1/f'(c) = 1/f'(g(d))$.
  3. Since $d$ was arbitrary, $g' = 1/(f' \circ g)$ on $J$.}
\end{remark*}

\begin{remark*}[Dependencies]
\begin{itemize}
  \item \hyperref[thm:inverse-function-theorem-monotone]{Inverse Function
    Theorem (Monotone Version)}: The pointwise result applied at each
    $d \in J$.
\end{itemize}

\FlashUses{thm:inverse-function-theorem-monotone}
\end{remark*}